% !TeX program = xelatex
% UTF-8. Compile from this directory or from the repository root with XeLaTeX.
% To input this file in an existing XePersian document, define \FWLSInNotebook.
% Source: Google Drive Notebook/FWLS, handwritten pages 1--9, filename order.
% pic01.jpg = 20261012_175032462.jpg; 1sHLKNnjPeRQgKOncd3kiHXPTmDB785hI
% pic02.jpg = 20261012_175036807.jpg; 1sw46zyw-XIZsQNFEDIMQQZC1LYH3IgRN
% pic03.jpg = 20261012_175040533.jpg; 1xu6XuSwO4f_2hsKF3DB4EF5Om6L8QBU9
% pic04.jpg = 20261012_175045713.jpg; 1ABg6M0rJPrgOoHyewHa1eZLiK0uoK7a1
% pic05.jpg = 20261012_175050754.jpg; 1WXOHLvK8QTsreagMZ9axh-JnfH_KIbod
% pic06.jpg = 20261012_175055957.jpg; 1hAtwOi2Ueyx_UjppHdKCK6JSt71mYwaO
% pic07.jpg = 20261012_175100732.jpg; 1n1zx0wS4Cmqhbvj_YnS_QPRRFjVej_-g
% pic08.jpg = 20261012_175107754.jpg; 1wercVZAOIsMtkGtBueMPs19zwAvgtEUc
% pic09.jpg = 20261012_175112659.jpg; 1MvJ7WfgsQ38g2ndQP9ICnzE73mYc4Xc6
\ifdefined\FWLSInNotebook
  \let\FWLSfinish\relax
\else
  \def\FWLSfinish{\end{document}}
  \documentclass[12pt,a4paper]{article}
  \usepackage[margin=22mm]{geometry}
  \usepackage{amsmath,amssymb}
  \usepackage{graphicx}
  \usepackage{xepersian}
  \IfFontExistsTF{Amiri}{\settextfont{Amiri}}{\settextfont{DejaVu Sans}}
  \setlatintextfont{DejaVu Serif}
  \setlength{\parindent}{0pt}
  \setlength{\parskip}{5pt}
  \setlength{\emergencystretch}{3em}
  \begin{document}
\fi
\providecommand{\FWLSunclear}[1]{\textbf{[خوانش نامطمئن: #1]}}
\providecommand{\FWLSmath}[1]{\lr{\(#1\)}}
\providecommand{\FWLSpage}[2]{\clearpage\section*{برگهٔ #1: \lr{\texttt{pic#2.jpg}}}}
\providecommand{\FWLSheader}[1]{\textbf{سربرگ و حاشیه:} ۲۲۰ ـ \lr{FWLS}؛
  \lr{Subject:} (خالی)؛ \lr{Date: / /}؛ شمارهٔ برگه: #1؛ نشان پایین برگه: \lr{YASHA}.\par}

\begin{center}
  {\Large\bfseries\lr{FWLS}}\\
  \lr{Nelson 22}\\
  ۲۲۰
\end{center}
\textbf{یادداشت رونویسی:} متن فارسی و انگلیسی، سربرگ‌ها، حاشیه‌ها و اصلاحات هر نه تصویر به ترتیب برگه‌ها آورده شده است. فلش‌ها در فهرست‌ها و روابط حفظ شده‌اند. خوانش‌های نامطمئن مشخص شده‌اند؛ گزاره‌های پزشکی منبع اصلاح نشده‌اند. در روابط عددیِ دارای جهت مبهم، ترتیب افقی نمادهای خود تصویر ثبت شده است. اصل نه تصویر در پیوست همین سند آمده است.

\FWLSpage{۱}{01}
\textbf{سربرگ و حاشیه:} ۲۲۰ ـ \lr{FWLS}؛ \lr{Nelson 22}؛
\lr{Subject:}؛ \lr{Date: / /}؛ شمارهٔ برگه: ۱؛ \lr{YASHA}.

\lr{FWLS} یا \lr{FWF}: دمای رکتال \FWLSmath{38^\circ\leq} بدون هیچ علائم و نشانهٔ بالینی دیگر.
\FWLSunclear{اختصار دوم به صورت FWF نوشته شده است. ترتیب عدد و علامت دما عین تصویر است.}

\subsection*{گروه‌های سنی در مواجهه با \lr{FWLS}}
\begin{itemize}
\item نوزادان: \lr{0--28} روز.
\item شیرخواران کوچک: \lr{29--90} روز.
\item کودکان: ۳ الی ۳۶ ماه.
\end{itemize}

\subsection*{اتیولوژی و اپیدمیولوژی}
\begin{itemize}
\item \lr{SBI} (عفونت باکتریال جدی) در \lr{7--13} درصد نوزادان و شیرخواران کوچک با تب دیده می‌شود.
\item شایع‌ترین علت \lr{UTI}: \lr{5--13} درصد.
\item باکتریمی: \lr{1--2\%}؛ مننژیت: \lr{0.2--0.5\%}.
\item ریسک \lr{SBI} در کودکان \lr{ill} و دارای ریسک فاکتور بالاتر است.
\item هرچه سن بالاتر، خطر \lr{SBI} کمتر.
\item \lr{IBI} (عفونت باکتریال مهاجم): باکتریمی و مننژیت.
\item باید حتماً \lr{ABT} تزریقی انجام دهیم (برخلاف مثلاً \lr{UTI} که می‌توانیم درمان خوراکی بدهیم).
\end{itemize}
\textbf{یادداشت اختصارات:} حرف آخر \lr{SBI} و \lr{IBI} در بعضی سطرها با خط افقی نوشته شده و شبیه \lr{SBT} و \lr{IBT} است؛ در سایر سطرها \lr{SBI} و \lr{IBI} دیده می‌شود.

\begin{itemize}
\item \lr{E. coli} شایع‌ترین ارگانیسم عامل \lr{SBI} است؛ سپس از آن \lr{GBS}.
\item سایر ارگانیسم‌ها: کلبسیلا، انتروکوک، استرپ پنومونه، نایسریا مننژیتیدیس، استاف اورئوس.
\item معمولاً فقط تب نمی‌دهند، در غالب پنومونی، \lr{GE}، \ldots
\item \lr{HSV} هم باید در نوزاد تب‌دار مدنظر باشد.
\item سن زیر ۲۸ روز.
\end{itemize}
\FWLSunclear{در جملهٔ «معمولاً فقط تب نمی‌دهند، در غالب پنومونی، GE، ...»، واژهٔ «غالب» همین‌گونه خوانده شده است.}

\FWLSpage{۲}{02}
\FWLSheader{۲}
\textbf{حاشیهٔ پایین راست:} عدد جداگانهٔ ۳ نیز کنار شمارهٔ زیرخط‌دار ۲ دیده می‌شود.

\subsection*{\lr{Neonatal HSV}}
\begin{itemize}
\item نادر است: \lr{0.2--0.3\%} نوزادان تب‌دار.
\item بیشتر در زمینهٔ تیپ \lr{II}.
\item درگیری به صورت \lr{SEM} عموماً در روزهای \lr{5--12} زندگی رخ می‌دهد.
\item توضیح \lr{SEM}: \lr{Skin / Eye / Mouth}.
\item درگیری به صورت درگیری \lr{CNS} عموماً در روز \lr{14--19} زندگی.
\item \lr{HSV} اکتسابی از مادر به صورت پره‌ناتال ممکن است پس از ۲۸ روزگی بروز کند.
\end{itemize}
\FWLSunclear{بازهٔ CNS در دست‌نویس «۱۹--۱۴» در راست‌به‌چپ خوانده می‌شود؛ واژهٔ پیش از «ممکن است» شبیه «پره‌ناتال» است.}

نوزاد تب‌دار \lr{appear well}: عفونت‌های ویروسی بسیار شایع‌تر از باکتریال.

شایع‌ترین ویروس: \lr{RSV}، انتروویروس، آنفلوآنزا، پاراآنفلوآنزا، \lr{h-MPV}، آدنوویروس.

\subsection*{تظاهرات بالینی}
در نوزادان و شیرخواران کوچک ممکن است عفونت‌های ویروسی و باکتریال صرفاً در قالب یک تب بروز کنند.

اما گاهی بیماران علائم مختلفی دارند:
\begin{itemize}
\item دمای غیرطبیعی: \FWLSmath{<36^\circ\ \mathrm{or}\ 38^\circ\leq}.
\item یافته‌های تنفسی: \FWLSmath{60<RR}، دیسترس تنفسی، آپنه.
\item یافته‌های همودینامیک: \FWLSmath{180<PR}، \FWLSmath{3<CR} ثانیه، نبض ضعیف.
\item یافته‌های معاینهٔ شکم.
\item یافته‌های غیرطبیعی نورولوژیک.
\item لتارژی، تحریک‌پذیری، تغییر تون.
\item یافته‌های غیرطبیعی پوست: راش، پتشی، سیانوز.
\end{itemize}
\FWLSunclear{علامت‌ها و ترتیب عدد و متغیر در سطرهای دما، RR، PR و CR مطابق تصویر ثبت شده‌اند. اختصار دست‌نویسِ زمان پرشدن مویرگی CR است.}

\FWLSpage{۳}{03}
\FWLSheader{۳}
\begin{itemize}
\item شیرخواران با آرتریت سپتیک یا استئومیلیت ممکن است علائم سپتیک نداشته باشند و یافته‌ها صرفاً محدود به ناحیهٔ درگیر شود.
\item یا یافتهٔ خودش را به صورت سودوپارالیز (\lr{disuse}) و یا یافتهٔ لوکال ایریتاسیون نشان دهد.
\item شیرخوار حین زمانی که والد تلاش می‌کند او را آرام کند، آرام می‌گیرد.
\end{itemize}
\FWLSunclear{در سطر دوم، عبارت پس از «یافته» شبیه «خودش را» است. در سطر سوم، واژهٔ پس از «شیرخوار» ممکن است «همین» باشد؛ فعل پایانی «آرام می‌گیرد» یا «آرام نمی‌گیرد» خوانده می‌شود.}

\subsection*{تشخیص}
قبلاً همهٔ بیماران \FWLSmath{<90} یا \FWLSmath{<60} روز را بستری می‌کردند ولی در حال حاضر بر اساس گایدلاین \lr{AAP} به \lr{FWLS} در شیرخوار \lr{0--36} ماه اپروچ می‌شود.

\begin{latin}
\subsection*{AAP}
\textbf{Any Toxic Appearing child: 0--36 months old + \(38^\circ\leq T\)}
\begin{itemize}
\item \(\longrightarrow\) Hospitalize; CBC; Culture (B/C, U/C, CSF) + Others (CXR, S/E, HSV or other virus PCR) + Parentral ABT.
\end{itemize}

\textbf{Well Appearing Child}
\begin{itemize}
\item Child \(<22\) days + \(38^\circ\leq T\)
  \(\longrightarrow\) Hospitalize, CBC, Culture (B/C, U/C, CSF) + Others + Parentral ABT.
\item Child 22--60 days + \(38^\circ\leq T\) \(\Longrightarrow\) 3 Step process.
\end{itemize}

\textbf{1. Determine Risk (H/x, Ph/E, Lab)}

\textbf{Low Risk}
\begin{itemize}
\item \(\longrightarrow\) Uncomplicated H/x.
\item \(\longrightarrow\) Well Appearing Ph/E.
\item \(\longrightarrow\) N/L Lab studies.
\item \(\longrightarrow\) U/A: Neg LE, Neg nitrite + \(\leq5\) WBC Centrifuged;
  \(<10\) WBC Uncentrifuged.
\end{itemize}
\end{latin}
\textbf{یادداشت املایی:} واژهٔ انگلیسی درمان تزریقی در دست‌نویس شبیه \lr{Parentral} نوشته شده و همین املا در رونویسی نگه داشته شده است.

\FWLSpage{۴}{04}
\FWLSheader{۴}
\begin{latin}
\subsection*{Low Risk (continued)}
\begin{itemize}
\item \(\longrightarrow\) Inf markers: \(T\leq38.5^\circ\), Procalcitonin \(\leq0.5\) ng/ml,
  CRP \(\leq20\) mg/L, ANC \(\leq4{,}000\text{--}5{,}200/\mathrm{mm}^3\).
\item \(\longrightarrow\) S/E if Diarrhea: No RBC and \(<5\) WBC/hpf.
\end{itemize}

\textbf{2. If child fulfills all low risk criteria}
\(\longrightarrow\) Use Age to determine need for LP, Parentral ABT and Hospitalization.
\begin{itemize}
\item Age 22--28: U/A, B/C, Inf markers, May perform LP, May administer parental ABT, Hospitalize.
\item Age 29--60: U/A, B/C, Inf markers, Need not perform LP,
  Need not administer ABT, Observe at home, follow-up 24--36.
\end{itemize}

\textbf{3. If child does not fulfill all low risk criteria}
\(\longrightarrow\) Use Age and Lab to determine need for LP, ABT, Hospitalization.
\begin{itemize}
\item Age 22--28 + ab N/L U/A + N/L inf markers
  \(\longrightarrow\) May perform LP, Administer parental ABT, Hospitalize.
\item Age 29--60 + ab N/L U/A + N/L inf markers
  \(\longrightarrow\) Oral ABT, Observe at home, follow-up 12--24 hr.
\item Age 22--28 + ab N/L inf markers
  \(\longrightarrow\) Perform LP, Parentral ABT only if (CSF pleocytosis, CSF uninterpretable, ab N/L U/A)
  and Hospitalize if parental ABT.
  \par If U/A and CSF \(\longrightarrow\) Observe at home after iv ABT.
\item Age 29--60 + ab N/L inf markers
  \(\longrightarrow\) May perform LP.
  \par If CSF pleocytosis \(\longrightarrow\) iv ABT and hospitalize.
  \par If CSF N/L \(\longrightarrow\) iv or PO ABT + observe home/hospital.
  \par If CSF not available or uninterpretable \(\longrightarrow\) iv ABT + observe at home.
\end{itemize}
\end{latin}
\textbf{یادداشت رونویسی برگهٔ ۴:} در سطر \lr{follow-up 24--36} واحدی نوشته نشده است. عبارت \lr{If U/A and CSF} نیز در منبع بدون صفت بعد از \lr{CSF} است. املاهای \lr{Parentral} و \lr{parental} هر دو در نوشته دیده می‌شوند.

\FWLSpage{۵}{05}
\FWLSheader{۵}
\textbf{حاشیهٔ بالا:} ۲۲۰ = \lr{FWLS}، افزون بر سربرگ ۲۲۰ ـ \lr{FWLS}.

\begin{latin}
\subsection*{Child 2--36 mo}
\textbf{Child 2--36 mo + \(38^\circ<T\leq39^\circ\)}
\begin{itemize}
\item \(\longrightarrow\) Reassure that diagnosis is likely self-limited viral infection,
  but advise return with persistence of fever, \(T>39\), and/or new signs and symptoms.
\end{itemize}

\textbf{Child 2--36 mo + \(39^\circ\leq T\)}
\(\longrightarrow\) Two-step Process.
\begin{enumerate}
\item Determine immunization status.
\item If received conjugate pneumococcal and Haemophilus influenzae type b vaccines
  \(\longrightarrow\) U/A, U/C for all females, all males \(<6\) mo,
  all uncircumcised males \(<2\) yr, and all children with recurrent UTI.
\end{enumerate}
\end{latin}
\FWLSunclear{عدد آغاز بازهٔ سنی در هر دو سطر انگلیسیِ این برگه شبیه 2 است و به صورت 2--36 حفظ شده؛ در برگهٔ ۱ گروه کودکان ۳ تا ۳۶ ماه نوشته شده است.}

\begin{itemize}
\item کودک \lr{ill}: باید ظرف ۱ ساعت احیا، بررسی و \lr{ABT} شود.
\item با توجه به ذکر این نکته، ریسک \lr{SBI} به دنبال مثبت شدن تست‌های ویروسی کاهش می‌یابد.
\item شیرخوار کوچک تب‌دار با برونشیولیت: ممکن است نیازمند \lr{LP} نباشد و بتوان به جای \lr{observe}، \ldots
\end{itemize}
\FWLSunclear{پایان جملهٔ برونشیولیت در منبع با سه‌نقطه رها شده است؛ عبارت «و بتوان به جای observe» خوانش محتمل است.}

پلئوسیتوز استریل مایع \lr{CSF} در ۲۳\% شیرخواران با \lr{UTI} دیده می‌شود.
\FWLSunclear{عدد این سطر «۲۳\%» خوانده می‌شود.}

عدد \lr{WBC} به تنهایی برای ارزیابی احتمال \lr{SBI} مناسب نیست.

۳۰\% شیرخواران با باکتریمی و ۴۰\% شیرخواران با مننژیت \lr{WBC} نرمال (\lr{5--15}) دارند.
\FWLSunclear{درصد دوم «۴۰\%» خوانده شده؛ صفرِ درصدها روی خط نوشته شده است. برای بازهٔ WBC در منبع واحد یا ضریب هزار ذکر نشده است.}

\FWLSpage{۶}{06}
\FWLSheader{۶}
\lr{B/C} منفی ردکنندهٔ احتمال مننژیت نیست.

بررسی \lr{CSF} شامل: کشت، رنگ‌آمیزی گرم (\lr{Stain})، \lr{cell count}، گلوکز، پروتئین، \lr{PCR} (بر اساس سناریو).

\begin{itemize}
\item \FWLSmath{\frac{\mathrm{WBC}}{\mathrm{RBC}}} در \lr{CSF}: \FWLSmath{\frac{1}{500}} در موارد تروماتیک؛ اصلاح فرمول.
\item \lr{WBC ratio}: \FWLSmath{\leq0.01}؛ رد مننژیت باکتریال.
\end{itemize}
\FWLSunclear{علامت نسبت WBC در تصویر پیش از 0.01 به صورت \FWLSmath{\leq} خوانده شده و ترتیب نمادها حفظ شده است.}

کشت برای نایسریا در \lr{CSF} ظرف ۲ ساعت پس از \lr{ABT} منفی می‌شود.

در مننژیت باکتریال پنوموکوک، \lr{CSF} ظرف \lr{4--24} ساعت بعد از \lr{ABT} به حد \lr{N/L} برمی‌گردد.

\textbf{اصلاح بالای سطر:} «استرپ پنومونه»؛ «۲ تا ۲۴».
\FWLSunclear{در عبارت اصلی «۴--۲۴» و در نوشتهٔ افزودهٔ بالای آن «۲ تا ۲۴» خوانده می‌شود؛ هر دو ثبت شده‌اند.}

پروتئین \lr{CSF} برای بیش از ۲۴ ساعت.

تغییرات \lr{WBC} و \lr{ANC} ظرف ۲۴ ساعت اول ناچیز است.

\subsection*{\lr{HSV}}
فقط فاکتور ریسک مشخصی وجود ندارد که کدام بیمار باید از نظر \lr{HSV} بررسی شود:
\FWLSunclear{آغاز سطر به صورت «فقط فاکتور ریسک مشخصی وجود ندارد» خوانده می‌شود.}
\begin{itemize}
\item شک: مواجهه با ۱ فرد بیمار، مثلاً مادر با سابقهٔ هرپس ژنیتال اولیه.
\item تشنج، یافتهٔ غیرطبیعی نورولوژیک.
\item راش وزیکولار.
\item ظاهر \lr{ill}، آپنه، هایپوترمی، پتشی یا خونریزی.
\item نارسایی کبد.
\item سابقهٔ آلکترود اسکالپ.
\end{itemize}
\FWLSunclear{واژهٔ پایانی سطر ظاهر ill به صورت «خونریزی» خوانده شده؛ عبارت آخر به صورت «سابقهٔ آلکترود اسکالپ» نوشته شده است.}

\FWLSpage{۷}{07}
\FWLSheader{۷}
\subsection*{اپروچ‌های موجود در خصوص درمان \lr{HSV} احتمالی}
\begin{enumerate}
\item تست و درمان امپریکال همهٔ نوزادان زیر ۲۱ روز که از نظر عفونت بررسی می‌شوند.
\item تست و درمان امپریکال نوزادانی که ویژگی‌های بالینی \lr{high-risk} برای \lr{HSV} دارند.
\item تست و درمان امپریکال همهٔ نوزادان با ویژگی‌های بالینی \lr{high-risk}
  + ارزیابی \lr{CSF} همهٔ نوزادان زیر ۲۱ روز بدون درمان امپریکال.
\end{enumerate}

نوزادی را که می‌خواهیم از نظر \lr{HSV} بررسی کنیم:
\begin{itemize}
\item کشت یا \lr{PCR}: سطح کنژنکتیوا، دهان، نازوفارنکس، رکتوم، هر ضایعهٔ وزیکولی.
\item \lr{CSF PCR}، \lr{Blood PCR}، \lr{ALT}.
\end{itemize}
\FWLSunclear{عبارت حاشیهٔ چپِ محل نمونه‌گیری به صورت «هر ضایعهٔ وزیکولی» خوانده شده است.}

بررسی انتروویروس: علت خوش‌خیم برای تب نوزادان؛ \lr{CSF PCR}.

\lr{PCT} در تشخیص \lr{IBI} در زیر ۳ سال از \lr{CRP} و \lr{WBC} بهتر است.

اگر سن \FWLSmath{21<} روز باشد، \lr{PCT} و \lr{CRP} منفی باشد، ظاهر \lr{ill} نباشد،
\FWLSmath{10{,}000>ANC} و پیوری وجود نداشته باشد: تقریباً رد کامل \lr{IBI}.

\textbf{اصلاح و حاشیه:} علامت سن ابتدا به صورت \FWLSmath{\leq21} نوشته و خط خورده؛ بالای آن \FWLSmath{21<} آمده است. در حاشیهٔ چپ \FWLSmath{10{,}000>} به \lr{ANC} متصل است؛ عدد در تصویر به صورت «۱۰،ـ» نوشته شده است.

در بیمار \FWLSmath{21<} روز، فاکتورهای التهابی باید بررسی شوند:
\begin{itemize}
\item یا \lr{PCT} + \lr{ANC / CRP}.
\item یا \lr{CRP} + \lr{ANC} + افزایش دما.
\end{itemize}
\textbf{یادداشت جهت نمادها:} دو عبارت سن \FWLSmath{21<} با همان ترتیب افقی تصویر حفظ شده‌اند.

\FWLSpage{۸}{08}
\FWLSheader{۸}
\lr{CXR} در بررسی شیرخوار تب‌دار تنها زمانی کاربرد دارد که علائم تنفسی وجود داشته باشد.

\subsection*{درمان \lr{ABT}}
\begin{itemize}
\item رژیم ۱: یک سفالوسپورین نسل ۳، معمولاً سفیم یا سفتازیدیم.
\item رژیم ۲: یک سفالوسپورین نسل ۳ + آمپی‌سیلین.
\item رژیم ۳: یک آمینوگلیکوزید + آمپی‌سیلین.
\end{itemize}
\FWLSunclear{نام داروی نخست در رژیم ۱ کوتاه و شبیه «سفیم» است؛ امکان خوانش «سفوتاکسیم» نیز وجود دارد و بدون قطعیت جایگزین نشده است. داروی دوم «سفتازیدیم» خوانده می‌شود.}

آمپی‌سیلین:
\begin{itemize}
\item درمان انتخابی \lr{GBS} است.
\item لیستریا مونوسایتوژنز را نیز پوشش می‌دهد.
\item اغلب انتروکوک‌ها و برخی \lr{E. coli}ها را پوشش می‌دهد.
\end{itemize}

برای نوزادان تا ۲۸ روز رژیم ۲ و ۳ توصیه شده است؛ چون خطر لیستریا وجود دارد.

برای شیرخوار کوچک‌تر (بالای ۲۸ روز)، رژیم ۱ (سفتریاکسون) توصیه می‌شود.
\FWLSunclear{صفت پس از «شیرخوار» در این سطر «کوچک‌تر» خوانده می‌شود.}

در موارد بیمار \lr{ill} و یا زمانی که \lr{gram-stain} مایع \lr{CSF} مثبت است: اضافه کردن ونکومایسین به رژیم قبلی.

نوزاد که قرار است آسیکلوویر بگیرد: \lr{60 mg/kg/day}، به دلیل درمان امپریکال.
\FWLSunclear{عبارت پایانی به صورت «به دلیل درمان امپریکال» خوانده شده است.}

\FWLSpage{۹}{09}
\FWLSheader{۹}
\subsection*{نوشته‌های بالای برگه و روی سربرگ}
\begin{itemize}
\item \lr{CSF}: کشت ظرف ۲۴، ۱۰۰\%.
\item \lr{B/C}: ظرف ۲۴ ساعت ۹۱\% مثبت و ۳۶ ساعت ۹۹\%.
\item عدد جداگانهٔ \lr{82} بالای ۱۰۰\% و عدد جداگانهٔ \lr{86} در سمت چپ بالای سربرگ.
\end{itemize}
\FWLSunclear{دو عدد کوچک بالای سربرگ به صورت 82 و 86 خوانده می‌شوند؛ جایگزین یا ارتباطشان با عددهای سطر پایین روشن نیست، بنابراین جداگانه ثبت شده‌اند.}

\subsection*{ترخیص از بیمارستان}
\begin{itemize}
\item نوزادی که حین بستری از نظر \lr{HSV} بررسی شده است نمی‌تواند ترخیص شود تا زمان آماده شدن جواب تست‌ها.
\item در بیماران طول بستری وابسته به نوع پاتوژن و محل عفونت است.
\end{itemize}

\clearpage
\section*{پیوست: نه تصویر کامل منبع}
\IfFileExists{images/pic01.jpg}{%
  \def\FWLSimagebase{images/}}{%
  \def\FWLSimagebase{pediatrics/ped-infectious diseases/FWLS/images/}}
\providecommand{\FWLSoriginal}[2]{%
  \clearpage\subsection*{اصل برگهٔ #1: \lr{\texttt{pic#2.jpg}}}%
  \begin{center}
  \includegraphics[width=\linewidth,height=0.87\textheight,keepaspectratio]{\FWLSimagebase pic#2.jpg}
  \end{center}}
\FWLSoriginal{۱}{01}
\FWLSoriginal{۲}{02}
\FWLSoriginal{۳}{03}
\FWLSoriginal{۴}{04}
\FWLSoriginal{۵}{05}
\FWLSoriginal{۶}{06}
\FWLSoriginal{۷}{07}
\FWLSoriginal{۸}{08}
\FWLSoriginal{۹}{09}
\FWLSfinish
